%! TEX program = xelatex
\documentclass[12pt]{article}

\usepackage[margin=1in]{geometry}
\usepackage{fontspec}
\usepackage{xeCJK}
\setCJKmainfont{SimSun}[AutoFakeBold=true]
\setmainfont{Times New Roman}
\usepackage{setspace}
\usepackage{amsmath,amssymb}
\usepackage{booktabs}
\usepackage{longtable}
\usepackage{tabularx}
\usepackage{array}
\usepackage{enumitem}
\usepackage{hyperref}
\usepackage[round,authoryear]{natbib}
\usepackage{xcolor}
\usepackage{tcolorbox}

\setstretch{1.15}
\setlength{\parindent}{0pt}
\setlength{\parskip}{0.6em}

\title{Early AI Errors and Human Delegation in Repeated Strategic Belief Tasks:\\An EWA-Inspired Experimental Study\\
\large Early AI Errors and Human Delegation in Repeated Strategic Belief Tasks\\
\large 早期 AI 错误与重复策略信念任务中的人类授权：基于 EWA 的实验研究}
\author{}
\date{}

\begin{document}

\maketitle

% ============== 摘要 ==============
\begin{abstract}
随着生成式人工智能进入预测、判断与辅助决策场景，人类面对 AI 的问题已经不只是是否信任某一次输出，而是如何在重复反馈中学习 AI 适合承担何种决策角色。既有 algorithm aversion、trust in automation 与 AI-assisted decision-making 文献大多关注单轮建议采纳、局部信任校准或短期错误反应，较少考察在多轮反馈中，AI 的早期错误是否会改变人类对 AI 决策角色的长期授权。

本研究提出，在重复策略信念任务中，即使 AI advisor 与 human analyst 在长期表现上完全等价，AI 的早期错误也可能比人类分析师的早期错误引发更强的负向授权更新。换言之，人类学习的并不只是 AI 是否准确，而是 AI 是否适合在决策系统中承担建议者、默认判断者或被授权代理人的角色。

为检验这一机制，本文设计一项 20 轮 repeated decision game with feedback。实验采用 varying-parameter beauty contest 作为基础任务。参与者每轮需要完成一个高阶信念预测任务，并在五种协作架构中选择一种：自我判断、AI 建议、人类分析师建议、AI 默认/委托、人类分析师默认/委托。每轮结束后，参与者观察自身收益、AI 预测误差和 human analyst 预测误差，并在下一轮重新选择协作架构。

实验采用四条件 between-subjects 设计：AI Early Error、Human Early Error、AI Late Error 与 AI Distributed Error。关键识别前提是，AI advisor 与 human analyst 在 20 轮中的胜率、平均绝对误差、错误严重程度分布和完全委托收益均被严格控制为等价。四组之间唯一系统差异是 focal errors 的来源与时点。由此，若 AI Early Error 组在后续轮次中持续降低 AI 授权，则该结果不能被解释为 AI 客观能力较差，而应被理解为人类对 AI 错误的非对称角色学习。

本文将 reduced-form 动态分析作为主证据，并使用 EWA-inspired attraction updating model 作为机制模型。五种人机协作方式被视为五个可选策略，每个策略具有随经验反馈更新的 attraction。EWA 模型用于检验 AI 错误是否比 human analyst 错误更强地降低 AI-related architectures 的吸引力。本文预期贡献在于，将 algorithm aversion 从单轮建议采纳问题推进为有限重复反馈环境中的授权学习问题，并为 AI 在高阶信念任务中的角色配置提供实验机制证据。

\vspace{1em}
\noindent\textbf{关键词：} AI delegation; algorithm aversion; repeated decision game; beauty contest; EWA; role learning; human-AI collaboration; asymmetric learning
\end{abstract}

% ============== 一、引言 ==============
\section{引言}

人工智能系统正在从信息工具转向预测、建议和自动决策系统。在这些场景中，人类不只是判断 AI 的某一次输出是否可信，而是需要在重复互动中决定 AI 应当被放置在决策流程中的什么位置。AI 可以只是提供参考信息，也可以给出建议，还可以形成默认判断，甚至在某些任务中被授权直接执行。因此，人类对 AI 的学习对象并非单纯的准确率，而是 AI 的角色适配性。

既有研究已经说明，人类对算法错误并非完全理性或 source-neutral。Algorithm aversion 研究表明，人们在观察到算法犯错后，可能会过度回避算法，即使算法总体表现并不差。Trust in automation 和 appropriate reliance 文献则强调，问题不在于人是否盲目信任系统，而在于人是否能够形成适当依赖。然而，这些研究大多关注单轮建议采纳、短期信任变化或局部依赖校准，较少解释人类如何在多轮反馈中学习 AI 的决策角色。

本文关注的核心问题是：在重复策略信念任务中，AI 的早期错误是否会导致人类在后续轮次中持续降低 AI 授权，即使 AI 与 human analyst 在长期表现上完全相同？

这一问题具有重要理论意义。如果 AI 和人类分析师长期胜率、平均误差和委托收益都相同，但参与者仍然因为 AI 早期错误而更少授权 AI，那么这种行为就不能被简单解释为对客观能力差异的理性反应。它更可能反映出一种非对称的角色学习机制：人类对 AI 错误赋予更高负向权重，并将早期 AI 错误解释为 AI 不适合承担更高层级的决策角色。

本文将这一机制放入一个 repeated decision game with feedback 中进行检验。参与者每轮面对一个高阶信念预测任务，并从多种人机协作架构中选择一种。每种协作架构都可以被看作一个可选策略。每轮结束后，参与者根据收益、AI 表现和 human analyst 表现更新对这些策略的主观吸引力，并在下一轮重新选择。该结构与 EWA，即 Experience-Weighted Attraction learning，的基本逻辑一致：策略选择并非一次性决定，而是在重复经验反馈中不断更新。

% ============== 二、理论框架 ==============
\section{理论框架}

\subsection{从信任 AI 到学习 AI 的角色}

传统 AI 信任研究通常关心人类是否相信 AI、是否采纳 AI 建议，或是否根据 AI 置信度调整依赖程度。但在真实的人机协作环境中，人类面对的往往不是单次建议采纳问题，而是授权结构选择问题。

例如，在同一类预测任务中，AI 可以扮演不同角色：
\begin{enumerate}[label=\arabic*.]
    \item 信息来源；
    \item 建议者；
    \item 默认判断者；
    \item 可被委托的代理人。
\end{enumerate}

这些角色在控制权、责任归因、复核成本和错误后果上存在差异。因此，本文将研究对象定义为 \textbf{role learning}，即个体在重复反馈中学习 AI 适合承担何种决策角色的过程。

\subsection{为什么选择 beauty contest}

本文采用 varying-parameter beauty contest 作为基础实验任务。该任务的核心特征是高阶信念推理：参与者不只是判断一个客观事实，而是预测他人会如何判断，以及他人如何预测其他人的判断。本文使用 beauty contest 并不是为了进入金融市场或投资决策研究，而是因为它提供了一个简洁、可控、可重复的策略信念任务，适合观察人在多轮反馈中如何更新对 AI 与 human analyst 的角色判断。

但是，标准 beauty contest 如果固定规则并重复过多轮，参与者可能很快学会"向均衡收敛"的套路。因此，本文不采用固定 2/3 average guessing game 进行 40 轮重复，而采用 \textbf{varying-parameter beauty contest}。具体而言，每轮的 $p$ 值、目标群体信息或历史信息呈现方式会发生变化，使任务保持同一理论结构，但避免机械收敛。

本研究选择 beauty contest 的原因有三点。

第一，它提供了一个基础、干净的高阶信念任务，而不是复杂的真实场景任务。\\
第二，它能够生成连续的预测误差，从而允许研究者严格控制 AI 与 human analyst 的表现等价。\\
第三，它适合被嵌入 repeated decision game with feedback，使参与者在多轮中更新对不同协作架构的偏好。

\subsection{Repeated decision game with feedback}

本文的实验并不需要被界定为标准两人 normal-form game。更准确的定义是：

\textbf{a repeated decision game with feedback}。

在每一轮中，参与者面对相同类型但参数变化的策略信念任务，并从五种协作架构中选择一种。每种架构都可以被理解为一个策略，决定预测权和决策控制权如何在参与者、AI advisor 与 human analyst 之间分配。

选择之后，参与者观察结果反馈，包括自己的收益、AI advisor 的预测误差和 human analyst 的预测误差。下一轮中，参与者根据此前经验重新选择协作架构。因此，本文的核心学习对象不是 beauty contest 中具体提交的数字，而是不同人机协作架构的 attraction。

\subsection{EWA 作为架构吸引力更新模型}

Experience-Weighted Attraction learning 的核心思想是，每个可选策略都有一个 attraction，该 attraction 会根据过去经验不断更新，并影响下一轮被选择的概率。本文将这一思想用于人机协作架构学习。

在本研究中，五种协作方式就是五个可选策略：
\begin{enumerate}[label=\arabic*.]
    \item Self-only；
    \item AI advice；
    \item Human analyst advice；
    \item AI default / AI delegation；
    \item Human analyst default / human delegation。
\end{enumerate}

每轮 beauty contest 产生预测误差和收益反馈。EWA-inspired model 则描述参与者如何将这些反馈转化为对不同架构的 attraction 更新。本文的核心结构命题是：AI 错误对 AI-related architectures 的负向更新权重大于 human analyst 错误对 human-related architectures 的负向更新权重。

% ============== 三、研究问题与假设 ==============
\section{研究问题与假设}

\subsection{研究问题}

本文的总研究问题是：

\begin{quote}
在重复策略信念任务中，当 AI advisor 与 human analyst 长期表现等价时，AI 的早期错误是否会比人类分析师的早期错误导致更强的后续低授权？
\end{quote}

具体而言，本文关注四个子问题：
\begin{enumerate}[label=\arabic*.]
    \item AI 早期错误是否降低后续 AI 授权？
    \item AI 早期错误的影响是否强于 human analyst 早期错误？
    \item AI 早期错误是否比 AI 后期错误和分散错误产生更持久影响？
    \item 这种行为轨迹是否可以通过 EWA-inspired attraction updating 解释？
\end{enumerate}

\subsection{研究假设}

\textbf{H1：AI 错误非对称惩罚假设。}\\
在 AI advisor 与 human analyst 长期表现等价的前提下，AI 早期错误将比 human analyst 早期错误引发更强的 AI 授权下降。

\textbf{H2：早期路径依赖假设。}\\
相较于 AI 后期错误，AI 早期错误将更强地降低后续 AI-related architectures 的选择概率。

\textbf{H3：集中错误冲击假设。}\\
相较于同等数量但分散出现的 AI 错误，早期集中 AI 错误将造成更强的后续低 AI 授权。

\textbf{H4：EWA 机制假设。}\\
在 EWA-inspired model 中，AI 错误的负向 learning weight 将显著大于 human analyst 错误的负向 learning weight，即：
\[
|\beta_{\text{AI}}^{-}| > |\beta_{\text{H}}^{-}|
\]

\textbf{H5：有限窗口中的持续性低授权假设。}\\
在 20 轮有限重复窗口中，AI Early Error 组将在后期轮次表现出持续性低 AI 授权。本文不将其表述为长期稳态锁定，而表述为 finite repeated-game horizon 中的 persistent post-error under-delegation。

% ============== 四、实验设计 ==============
\section{实验设计}

\subsection{总体设计}

本研究采用 20 轮在线行为实验。20 轮设计基于实验经济学中 repeated-game learning 和 belief formation 研究的常见做法。EWA 经典研究并不要求所有任务都具有超长轮次；在部分 normal-form games 中，10 轮数据已经被用于观察经验学习。关于 repeated games 中 learning、teaching 和 convergence 的研究也常采用 20 periods 设计。因此，20 轮能够在避免疲劳和机械收敛的同时，为观察授权更新和有限收敛提供足够窗口。

本文不以估计长期均衡或精细个体学习参数为主要目标，而以识别早期错误后的授权调整和有限重复窗口中的持续性低授权为目标。因此，20 轮比 40 轮更适合本研究。40 轮在 beauty contest 环境下可能导致任务套路化和疲劳；8–12 轮则不足以同时覆盖早期错误、后续调整、后期错误对照与恢复窗口。

\subsection{轮次安排}

20 轮被划分为五个阶段：

\begin{center}
\begin{tabular}{lcl}
\toprule
\textbf{阶段} & \textbf{轮次} & \textbf{功能} \\
\midrule
Baseline & 1--2 & 熟悉任务，观察 AI 与 human analyst 均正常表现 \\
Early shock & 3--5 & 早期 focal errors 出现 \\
Adjustment & 6--12 & 观察错误后的授权调整 \\
Late shock / balancing & 13--15 & AI Late 组出现 focal errors，其余组进行表现配平 \\
Recovery / finite convergence & 16--20 & 观察后期授权是否恢复或保持低授权倾向 \\
\bottomrule
\end{tabular}
\end{center}

该结构允许研究者同时观察 early shock、post-shock adjustment、late shock comparison 和短期恢复窗口。

\subsection{实验任务}

每轮中，参与者完成一个 varying-parameter beauty contest。参与者需要在 0 到 100 之间提交一个数字，目标是尽可能接近该轮的 winning target：

\[
\text{Target}_t = p_t \times \text{GroupAverage}_t
\]

其中 $p_t$ 在不同轮次变化，例如 0.45、0.55、0.65、0.75。$\text{GroupAverage}_t$ 可来自预先收集的独立样本或实验池中的真实群体选择。每轮结束后，参与者看到 $\text{Target}_t$，并根据提交值与 $\text{Target}_t$ 的距离获得收益。

收益函数为：
\[
\text{Payoff}_{it} = B - \alpha \times |\text{SubmittedNumber}_{it} - \text{Target}_t| - \text{Cost}_k
\]

其中 $B$ 为基础收益，$\alpha$ 为误差惩罚系数，$\text{Cost}_k$ 为协作架构 $k$ 的成本。

\subsection{五种协作架构}

每轮正式提交前，参与者从以下五种协作架构中选择一种：
\begin{enumerate}[label=\arabic*.]
    \item \textbf{Self-only}：完全由参与者自己判断并提交；
    \item \textbf{AI advice}：AI advisor 给出预测建议，参与者最终决定；
    \item \textbf{Human analyst advice}：human analyst 给出预测建议，参与者最终决定；
    \item \textbf{AI default / AI delegation}：AI 生成默认预测，参与者可接受或少量修改；在高授权版本中可直接委托 AI 提交；
    \item \textbf{Human analyst default / human delegation}：human analyst 生成默认预测，参与者可接受或少量修改；在高授权版本中可直接委托 human analyst 提交。
\end{enumerate}

五种架构构成 repeated decision game 中的五个可选策略。本文的主因变量不是参与者每轮具体猜测的数字，而是参与者在 20 轮中如何调整这些协作架构的选择。

\subsection{反馈信息}

每轮结束后，参与者观察以下信息：
\begin{enumerate}[label=\arabic*.]
    \item 自己的最终提交值；
    \item 该轮 $\text{Target}_t$；
    \item 自己的收益；
    \item AI advisor 的预测值与预测误差；
    \item Human analyst 的预测值与预测误差；
    \item 如果完全委托 AI 或 human analyst，本轮会获得的收益。
\end{enumerate}

显示 AI 与 human analyst 的预测误差和完全委托收益有两个目的。第一，它保证参与者即使没有选择 AI delegation，也能观察 AI 本轮表现，从而避免 treatment exposure 内生性。第二，它使 EWA-inspired model 可以允许未选择策略发生部分 counterfactual updating，而不只是纯 reinforcement learning。

% ============== 五、表现等价控制 ==============
\section{表现等价控制}

本文的关键识别前提是 AI advisor 与 human analyst 的长期表现等价。实验不比较 AI 与人类分析师谁更强，而是检验当二者能力相同但错误来源与时点不同的时候，人类是否仍会对 AI 进行更强负向授权更新。

因此，正式实验前将预注册完整的 20 轮 AI/human 预测序列，并确保以下四项等价条件。

\subsection{胜率等价}

定义命中变量：
\begin{align*}
\text{Hit}_{\text{AI},t} &= \mathbf{1}(|F_{\text{AI},t} - \text{Target}_t| \leq \tau) \\
\text{Hit}_{\text{H},t} &= \mathbf{1}(|F_{\text{H},t} - \text{Target}_t| \leq \tau)
\end{align*}

其中 $\tau$ 为预注册命中阈值。要求：
\[
\sum \text{Hit}_{\text{AI},t} = \sum \text{Hit}_{\text{H},t}
\]

\subsection{平均绝对误差等价}

要求：
\[
\text{Mean}|F_{\text{AI},t} - \text{Target}_t| = \text{Mean}|F_{\text{H},t} - \text{Target}_t|
\]

\subsection{错误严重程度等价}

将错误按照绝对误差划分为 low、medium、high 三类。例如：
\begin{align*}
\text{low error:} & \quad \tau < \text{error} \leq a \\
\text{medium error:} & \quad a < \text{error} \leq b \\
\text{high error:} & \quad \text{error} > b
\end{align*}

要求 AI 与 human analyst 在 low、medium、high error 的数量上完全一致。

\subsection{完全委托收益等价}

如果参与者 20 轮全部委托给 AI 或全部委托给 human analyst，二者累计收益必须相同或落在预注册容忍范围内：
\[
\sum \text{Payoff}_{\text{AI-delegation},t} = \sum \text{Payoff}_{\text{H-delegation},t}
\]

这四项等价控制是本文因果识别的基础。只有在 AI 和 human analyst 客观表现等价时，后续授权差异才可以被解释为错误来源和错误时点导致的非对称学习，而不是对真实能力差异的理性反应。

% ============== 六、实验条件 ==============
\section{实验条件}

本研究采用四条件 between-subjects 设计。

\begin{center}
\begin{tabular}{llll}
\toprule
\textbf{条件} & \textbf{focal error 来源} & \textbf{focal error 时点} & \textbf{识别目的} \\
\midrule
AI Early Error & AI advisor & Rounds 3--5 & 检验 AI 早期错误是否造成后续低授权 \\
Human Early Error & Human analyst & Rounds 3--5 & 检验 AI 错误与人类错误是否存在来源非对称 \\
AI Late Error & AI advisor & Rounds 13--15 & 检验早期错误是否比后期错误更具路径依赖 \\
AI Distributed Error & AI advisor & 分散于多个阶段 & 检验早期集中错误是否不同于分散错误 \\
\bottomrule
\end{tabular}
\end{center}

所有条件共享相同的 $\text{Target}_t$ 序列、$p_t$ 序列、收益函数和架构成本。四组之间唯一系统差异是 focal errors 的来源与时点。为了维持表现等价，非 focal errors 通过 balancing slots 安排，使 AI 与 human analyst 在 20 轮结束时满足胜率、平均误差、错误严重程度和完全委托收益的等价约束。

三组核心比较如下：
\begin{enumerate}[label=\arabic*.]
    \item \textbf{AI Early Error vs Human Early Error}：识别 source asymmetry；
    \item \textbf{AI Early Error vs AI Late Error}：识别 timing effect 和 early path dependence；
    \item \textbf{AI Early Error vs AI Distributed Error}：识别 early clustered shock effect。
\end{enumerate}

% ============== 七、样本计划 ==============
\section{样本计划}

正式实验计划招募 $N = 320$–$400$ 名参与者，每组约 80–100 人。最终样本量将根据 pilot study 的效应量估计和 power analysis 决定。

Pilot study 主要用于检验：
\begin{enumerate}[label=\arabic*.]
    \item 参与者是否理解 varying-parameter beauty contest；
    \item 五种协作架构是否清晰；
    \item 20 轮是否存在显著疲劳；
    \item focal errors 是否足够明显；
    \item AI 与 human analyst 表现等价是否能够被参与者感知；
    \item EWA-inspired 模型中的关键参数是否具有可识别性。
\end{enumerate}

正式实验将设置理解测试。未能理解任务规则、收益结构或协作架构含义的参与者不得进入正式付费轮次。

% ============== 八、变量与测量 ==============
\section{变量与测量}

\subsection{主因变量}
\begin{enumerate}[label=\arabic*.]
    \item \textbf{AI 授权比例}：选择 AI advice 或 AI default/delegation 的轮次比例；
    \item \textbf{高 AI 授权比例}：选择 AI default/delegation 的轮次比例；
    \item \textbf{后期 AI 授权比例}：Rounds 16--20 中 AI-related architectures 的选择比例；
    \item \textbf{低 AI 授权倾向}：后期持续选择 Self-only 或 human-related architectures；
    \item \textbf{架构切换率}：focal error 后从 AI-related architectures 转出的概率；
    \item \textbf{收益 benchmark 偏离}：参与者所选架构收益相对于预注册最优架构收益的偏离。
\end{enumerate}

\subsection{次级机制变量}

实验前后测量以下变量：
\begin{enumerate}[label=\arabic*.]
    \item 基线 AI 态度；
    \item 算法厌恶倾向；
    \item 风险态度；
    \item 控制偏好；
    \item 对 AI advisor 的能力评价；
    \item 对 human analyst 的能力评价；
    \item 对 AI 适合作为 advice/default/delegation 角色的评价；
    \item 对 human analyst 适合作为 advice/default/delegation 角色的评价；
    \item 错误归因：偶然波动、能力不足、角色不适配。
\end{enumerate}

这些变量用于区分 role learning、一般损失厌恶、控制偏好和 anti-AI prior。

% ============== 九、分析策略 ==============
\section{分析策略}

\subsection{Reduced-form 动态分析}

主分析采用 reduced-form dynamic treatment effect。基本模型为：
\begin{align*}
\text{AI\_Delegation}_{it} =\ & \beta_0 + \beta_1 \text{AIEarly}_i \times \text{PostShock}_t \\
& + \beta_2 \text{HumanEarly}_i \times \text{PostShock}_t \\
& + \beta_3 \text{AILate}_i \times \text{LateWindow}_t \\
& + \beta_4 \text{AIDistributed}_i \times \text{PostShock}_t \\
& + \gamma X_i + \delta_t + \varepsilon_{it}
\end{align*}

其中 $\text{AI\_Delegation}_{it}$ 表示参与者 $i$ 在 $t$ 轮的 AI 授权程度；$X_i$ 为个体控制变量；$\delta_t$ 为轮次固定效应。

核心检验包括：
\begin{enumerate}[label=\arabic*.]
    \item $\beta_1$ 是否显著为负；
    \item $\beta_1$ 是否显著小于 $\beta_2$；
    \item AI Early Error 组在 Rounds 16--20 是否仍保持较低 AI 授权；
    \item AI Early Error 的影响是否强于 AI Late Error 和 AI Distributed Error。
\end{enumerate}

\subsection{EWA-inspired 机制模型}

设 $k$ 表示五种协作架构。每个参与者 $i$ 对每个架构 $k$ 在轮次 $t$ 有一个 attraction：$A_{i,k,t}$。

更新方程为：
\[
A_{i,k,t} = \varphi A_{i,k,t-1} + \pi_{i,k,t}
\]

其中 $\varphi$ 表示经验保留程度。$\pi_{i,k,t}$ 表示该轮反馈：
\begin{align*}
\pi_{i,k,t} =\ & \theta \text{Payoff}_{i,k,t} \\
& + \mathbb{I}(k \text{ uses AI}) \times [\beta_{\text{AI}}^{+} \text{AI\_Correct}_t + \beta_{\text{AI}}^{-} \text{AI\_Error}_t] \\
& + \mathbb{I}(k \text{ uses Human}) \times [\beta_{\text{H}}^{+} \text{H\_Correct}_t + \beta_{\text{H}}^{-} \text{H\_Error}_t] \\
& - c_k
\end{align*}

下一轮选择架构 $k$ 的概率为：
\[
P(s_{i,t+1} = k) = \frac{\exp(\lambda A_{i,k,t} - \kappa \mathbb{I}(k \neq s_{i,t}))}{\sum_j \exp(\lambda A_{i,j,t} - \kappa \mathbb{I}(j \neq s_{i,t}))}
\]

其中 $\lambda$ 表示选择对 attraction 的敏感度，$\kappa$ 表示 switching cost。

核心机制检验为：
\[
|\beta_{\text{AI}}^{-}| > |\beta_{\text{H}}^{-}|
\]

若该条件成立，说明 AI 错误比 human analyst 错误更强地降低相关架构的 attraction。

由于本文只有 20 轮，EWA-inspired 模型将作为机制分析而非唯一主证据。主结论依赖 reduced-form 条件比较，EWA 用于解释观察到的授权轨迹是否符合非对称 attraction updating。

\subsection{模型比较}

本文将比较三个模型：
\begin{enumerate}[label=\arabic*.]
    \item \textbf{EWA-inspired source-asymmetric model}；
    \item \textbf{Source-neutral learning model}：AI 错误和 human 错误具有相同权重；
    \item \textbf{Inertia-only model}：选择主要由上一轮选择和 switching cost 决定。
\end{enumerate}

如果 source-asymmetric EWA-inspired model 显著优于后两者，并且 $\beta_{\text{AI}}^{-}$ 的负向权重大于 $\beta_{\text{H}}^{-}$，则支持本文的机制解释。

% ============== 十、效度与替代解释 ==============
\section{效度与替代解释}

\subsection{内部有效性}

本文的内部有效性依赖于表现等价控制和共同反馈设计。AI 与 human analyst 在胜率、平均误差、错误严重程度和委托收益上完全等价，排除了"AI 真的更差"的替代解释。每轮展示 AI 与 human analyst 的预测误差和假设委托收益，排除了参与者没有观察到某一来源表现的 exposure 内生性。

\subsection{构念有效性}

本文的核心构念是 role learning，而非单次 trust。为此，研究同时测量实际授权行为和 role-specific belief。如果参与者降低 AI 授权，同时降低对 AI 作为 default setter 或 delegated agent 的适配性评价，则更能支持角色学习解释。

\subsection{替代解释}
\begin{enumerate}[label=\arabic*.]
    \item 若结果来自风险厌恶或损失厌恶，则 AI Early Error 和 Human Early Error 应产生类似效果。
    \item 若结果来自一般行为惯性，则 inertia-only model 应足以解释选择轨迹。
    \item 若结果来自 anti-AI prior，则控制基线 AI 态度后处理效应应显著减弱。
    \item 若结果来自 AI 客观表现较差，则表现等价控制将直接排除该解释。
\end{enumerate}

% ============== 十一、预期贡献 ==============
\section{预期贡献}

\subsection{理论贡献}

本文将 algorithm aversion 从单轮建议采纳推进为有限重复反馈中的授权学习问题。研究不只问人是否在一次错误后回避 AI，而是问早期 AI 错误是否会改变人类对 AI 决策角色的持续授权。

\subsection{方法贡献}

本文提出一种将 varying-parameter beauty contest、repeated decision game with feedback、表现等价控制和 EWA-inspired 模型结合的实验框架。该框架能够在不依赖复杂真实业务场景的情况下，研究高阶信念环境中的 AI 角色学习。

\subsection{任务设计贡献}

本文将 beauty contest 从一个经典高阶信念任务改造为人机协作授权实验的反馈环境。其价值不在于模拟某个具体应用领域，而在于提供一个足够抽象、可控且可参数化的任务结构，使研究者能够严格控制 AI 与 human analyst 的表现等价，并检验错误来源与错误时点如何影响后续授权。

\subsection{AI 治理贡献}

如果 AI 早期错误在长期表现等价条件下仍导致持续性低授权，则说明 AI 部署初期具有特殊治理意义。组织不应只关注 AI 平均准确率，还应关注早期错误如何改变后续授权轨迹、恢复机会和用户角色认知。

% ============== 十二、结论 ==============
\section{结论}

本研究提出一项 20 轮 repeated decision game with feedback，用以检验 AI 早期错误是否会在高阶信念任务中造成持续性低授权。实验采用 varying-parameter beauty contest 作为反馈环境，将五种人机协作方式视为五个可选策略，并用 EWA-inspired model 描述参与者如何根据经验反馈更新不同架构的 attraction。

研究的核心识别原则是：AI advisor 与 human analyst 的总体表现完全等价，实验只操控错误来源与错误时点。由此，若 AI Early Error 组在后续轮次中持续降低 AI 授权，则说明人类对 AI 错误存在非对称角色学习，而不是仅仅根据客观表现进行理性调整。

本文最终试图建立一条清晰机制链：

\begin{center}
\fbox{\parbox{0.9\linewidth}{\centering
高阶信念预测任务 $\rightarrow$ AI/human 表现反馈 $\rightarrow$ 协作架构 attraction 更新 $\rightarrow$ 后续授权选择 $\rightarrow$ 有限重复窗口中的持续性低授权。
}}
\end{center}

% ============== 参考文献 ==============
\bibliographystyle{apalike}
\bibliography{references}

% ============== 问题备忘录 ==============
\appendix
\section{问题备忘录}

\begin{tcolorbox}[colback=blue!5!white, colframe=blue!75!black, title=\textbf{M1：20 轮到底够不够？}]
20 轮比 40 轮更可行，也比 8–12 轮更能覆盖学习过程；但它是否足以支持 "persistent post-error under-delegation" 仍然需要谨慎表述。现在最好不要说长期稳态，只说有限重复窗口中的持续性低授权。RP 现在已经采用这个更稳的说法。
\end{tcolorbox}

\begin{tcolorbox}[colback=blue!5!white, colframe=blue!75!black, title=\textbf{M2：选美博弈会不会太快被学会？}]
这是第二个核心风险。标准 beauty contest 重复多轮容易收敛，所以现在改成 varying-parameter beauty contest 是合理的。但还要进一步说明：每轮变化的是 $p_t$、历史信息或目标群体信息，而不是完全换任务。否则任务可能变得不统一。
\end{tcolorbox}

\begin{tcolorbox}[colback=blue!5!white, colframe=blue!75!black, title=\textbf{M3：四组条件在 20 轮内会不会太挤？}]
现在四组包括 AI Early、Human Early、AI Late、AI Distributed。理论上很完整，但 20 轮里塞 early shock、adjustment、late shock、recovery，确实紧。尤其 AI Late Error 之后只有 16–20 这 5 轮观察窗口，所以 AI Late 组最好作为辅助比较，而不是最核心比较。
\end{tcolorbox}

\begin{tcolorbox}[colback=red!5!white, colframe=red!75!black, title=\textbf{M4：AI 与 human analyst 的表现等价能不能真的做到？}]
这是实验识别的命门。RP 现在已经写了胜率、平均绝对误差、错误严重程度和完全委托收益四重等价控制，这很好。争议在于：20 轮很短，配平空间有限，所以必须做一个轮次矩阵，否则读者会怀疑等价控制只是原则声明。
\end{tcolorbox}

\begin{tcolorbox}[colback=red!5!white, colframe=red!75!black, title=\textbf{M5：EWA 会不会在 20 轮里估不稳？}]
这是方法上最大的风险。20 轮可以用 EWA-inspired model，但不要把它当强结构估计。主证据应该是 reduced-form 条件比较，EWA 只作为机制辅助，重点检验 $|\beta_{\text{AI}}^{-}| > |\beta_{\text{H}}^{-}|$，不要估太多复杂参数。
\end{tcolorbox}

\begin{tcolorbox}[colback=orange!5!white, colframe=orange!75!black, title=\textbf{M6：五种协作方式是否足够表达"角色学习"？}]
现在五种方式是 Self-only、AI advice、Human advice、AI delegation、Human delegation。它比原来更清楚，但理论范围也变窄了。严格说，它研究的是 advice vs delegation 的授权学习，而不是完整的人机协作架构演化。这里需要把理论表述收紧。
\end{tcolorbox}

\begin{tcolorbox}[colback=orange!5!white, colframe=orange!75!black, title=\textbf{M7：反馈信息给得太多，会不会削弱路径依赖？}]
现在每轮显示 AI/human 的预测误差和完全委托收益。好处是排除了"参与者没看到 AI 后续表现"的问题；坏处是路径依赖会更难出现，因为参与者有充分反事实信息。如果这种情况下仍出现 AI 低授权，结果会很强；但效应可能更小。
\end{tcolorbox}

\begin{tcolorbox}[colback=green!5!white, colframe=green!50!black, title=\textbf{M8：行为金融学连接够不够？}]
beauty contest 和高阶信念确实能接到行为金融学，但现在还需要避免被看成一般 HCI 实验。比较稳的说法是：本文不研究投资策略，而研究金融市场中更基础的 higher-order belief formation 与预测主体授权。
\end{tcolorbox}

\vspace{1em}
\noindent\textbf{备忘录图例：}
\begin{itemize}
    \item \textcolor{blue}{\textbf{蓝色}}：理论和方法框架已稳定，可在 RP 中正面陈述
    \item \textcolor{red}{\textbf{红色}}：识别与统计风险，需在 RP 中显性处理
    \item \textcolor{orange}{\textbf{橙色}}：构念范围与外部效度，需要进一步约束
    \item \textcolor{green!50!black}{\textbf{绿色}}：与外部学科连接，需重新表述
\end{itemize}

\end{document}